\section{O painel de dados e como ler este documento}
\label{sec:painel}

Esta parte do documento apresenta a Análise Exploratória dos dados da
pesquisa: antes de qualquer conclusão sobre o efeito do modelo de
gestão, é preciso conhecer em detalhe o material com que trabalhamos.
O que se faz aqui é responder perguntas simples e fundamentais: como se
distribuem os indicadores de desempenho entre os hospitais? Quem são os
hospitais de cada categoria de gestão e eles são comparáveis entre si?
O que aconteceu nos anos da pandemia? O que os dados mostram nos cinco
hospitais que trocaram a administração Direta pela gestão por OSS? E
onde estão os problemas de qualidade de registro que precisam de
cuidado?

\subsection{O painel}

A base é o painel consolidado do projeto, com \textbf{314 hospitais
acompanhados ano a ano de 2015 a 2025, totalizando 3.454 observações}
(uma observação corresponde a um hospital em um ano). O painel é
\textbf{balanceado}: todos os 314 hospitais aparecem nos 11 anos, sem
buracos, o que permite comparar a evolução de cada um deles em bases
iguais. Esse conjunto resulta de um funil documentado de seleção que
partiu de 830 estabelecimentos da base bruta do SIH/SUS e reteve
apenas hospitais com porte mínimo (mais de 50 leitos na mediana dos
anos), produção contínua no período e classificação de complexidade
disponível, excluindo ainda os hospitais de campanha da pandemia.

Os indicadores analisados são os cinco oficiais do estudo, sempre nas
versões que excluem as internações por COVID: \textbf{mortalidade
geral} (proporção de internações que terminam em óbito);
\textbf{mortalidade ajustada}, que exclui os óbitos fetais, maternos e
neonatais; \textbf{tempo médio de permanência} (TMP, em dias);
\textbf{custo por saída} (valor pago por internação encerrada, em
reais); e \textbf{fração de alta complexidade} (proporção das
internações classificadas como de alta complexidade). A eles se somam
as duas \textbf{taxas de ocupação} (internação e UTI), em escala
percentual.

Três regras valem para todo o documento. Primeira: a categoria de
gestão de cada hospital (Direta, OSS, Público Municipal, Filantrópico,
Privado) segue a classificação assistencial do SIH como definição
adotada pela equipe; a categoria \textbf{Privado reúne apenas 3
hospitais} e por isso nunca deve ser lida como retrato de um grupo,
apenas registrada. Segunda: o \textbf{porte} de cada hospital é fixo,
calculado sobre a mediana de leitos de todos os anos com os cortes
oficiais (até 50 leitos, pequeno porte; 51 a 150, médio; 151 a 500,
grande; acima de 500, especial); como o funil já excluiu os pequenos, o
painel tem \textbf{173 hospitais médios, 130 grandes e 11 especiais}, e
o rótulo anual do SIH, que oscila com o cadastro, não é usado.
Terceira: nas análises que envolvem mortalidade, a complexidade do
hospital entra sempre pela versão \textbf{estrutural} do escore
(baseada só em estrutura física e tecnológica), nunca pela versão que
usa a própria mortalidade, para não misturar causa e consequência.

\subsection{Guia rápido de leitura dos gráficos}

Alguns tipos de gráfico e de medida se repetem ao longo do documento;
vale fixá los uma única vez, em linguagem direta.

\textbf{Mediana e média.} A mediana é o valor do hospital do meio da
fila: metade dos hospitais está abaixo dela, metade acima. A média é a
soma dividida pelo total e é sensível a valores extremos: um único
hospital muito caro puxa a média para cima, mas quase não move a
mediana. Como nossos indicadores têm valores extremos legítimos, a
mediana é quase sempre o retrato mais fiel do hospital típico, e por
isso é a medida preferida aqui. Os \textbf{percentis} generalizam a
ideia: o percentil 90 (p90) é o valor abaixo do qual estão 90\% dos
hospitais.

\textbf{Histograma.} É a fotografia da distribuição: o eixo horizontal
traz o valor do indicador e a altura de cada barra mostra quantas
observações caem naquele trecho. Onde as barras são altas, os hospitais
se concentram; caudas compridas para a direita indicam poucos casos com
valores muito altos.

\textbf{Gráfico de violino.} É um histograma de pé e espelhado, um por
categoria: a silhueta engorda onde há mais hospitais. As linhas
internas marcam os quartis e o número impresso é a mediana do grupo.

\textbf{Gráfico QQ.} Compara os dados reais com um molde matemático de
distribuição. Se os pontos seguem a reta, o molde descreve bem os
dados; desvios sistemáticos da reta indicam onde o molde falha. O
número r que acompanha cada QQ resume essa aderência: quanto mais
próximo de 1, melhor.

\textbf{AIC.} É uma nota de qualidade de ajuste usada para comparar
moldes de distribuição entre si: \textbf{quanto menor o AIC, melhor o
molde descreve os dados}. A nota só serve para comparação; o valor
isolado não tem interpretação.

\textbf{Escala logarítmica.} Alguns gráficos usam eixo em escala
logarítmica, em que as marcas avançam por multiplicação (100, 1.000,
10.000) em vez de soma. Isso permite enxergar ao mesmo tempo hospitais
de R\$~500 e de R\$~20.000 por saída sem que os pequenos virem um
borrão; os rótulos dos eixos mostram sempre os valores reais.
